\documentclass[UTF8,11pt,a4paper]{ctexart}
\usepackage{geometry,booktabs,longtable,tabularx,array,fancyhdr}
\geometry{left=2.5cm,right=2.5cm,top=2.5cm,bottom=2.5cm}
\pagestyle{fancy}\fancyhf{}\cfoot{\thepage}\renewcommand{\headrulewidth}{0pt}
\setlength{\parindent}{2em}
\newcolumntype{Y}{>{\raggedright\arraybackslash}X}
\title{AI工具使用详情}
\author{}
\date{}
\begin{document}
\maketitle

\section{工具与使用范围}

本参考项目使用 OpenAI Codex 桌面环境中的 GPT-5 系列模型辅助完成资料整理、建模草案检查、Python代码实现与审查、实验解释、图表规划、LaTeX排版和文稿校对。应用可能对具体模型路由和推理配置进行管理，因此本记录不虚构不可见的精确内部版本号。另使用同一环境中的独立AI审查任务，对模型、程序输出和论文证据进行只读复核。

AI没有代替参赛队伍作最终学术判断，也没有执行登录、上传或提交。当前文稿属于参考初稿，尚不能声明已由队伍完成逐项人工核验。

\section{各阶段用途、主要输入与采纳情况}

\begin{longtable}{p{2.2cm}p{4.2cm}p{4.6cm}p{3.0cm}}
\toprule
阶段 & AI辅助目的 & 主要输入或操作 & 当前采纳状态\\
\midrule
\endfirsthead
\toprule
阶段 & AI辅助目的 & 主要输入或操作 & 当前采纳状态\\
\midrule
\endhead
题目理解 & 拆分四个子问题，识别10分钟时间口径、储能约束、合同调账和变价信息边界 & 读取题目PDF、五类附件字段与官方格式规则；形成变量、假设和风险点 & 已形成参考分析；需队伍对照原题复核\\
模型建立 & 构造统一母线平衡和SOC递推；设计相似日预测、联合残差聚类、边际分位轨迹和滚动LP & 讨论因果历史窗口、冷启动、经验分位数、终端价值、实际执行优先级 & 已用于代码和论文；参数选择需人工判断\\
程序实现 & 编写数据读取、线性规划、全年连续模拟、日内更新、结算、敏感性和结果表程序 & 依据附件格式实现Python；运行测试、334日统计和54组实验 & 数值已实际运行；仍需队伍阅读代码\\
图表与结果 & 选择原始数据、建模过程和结果图；检查单位、统计期、灰度可读性与解释边界 & 从CSV生成PNG/SVG；比较固定价、变价、更新频率和保守参数 & 图已进入参考稿；需人工核对图注\\
文献 & 搜索微网随机优化、MPC和滚动时域相关论文，核验书目信息 & 使用双源学术检索发现文献，并打开DOI元数据页面核对作者、题名、年份和来源 & 仅作为研究背景；未把元数据核验冒充全文阅读\\
论文写作 & 建立结论到公式、表格、图和代码输出的对应关系，生成中文LaTeX参考稿 & 使用真实结果撰写摘要、四题模型、验证、结论、AI声明和附录 & 属AI辅助参考文本；必须人工改写与确认\\
排版检查 & 编译PDF并检查页型、正文页数、公式、图表、引用、字体、空白页和图片分辨率 & XeLaTeX、BibTeX、PDF渲染及视觉抽检 & 技术检查已完成；不等于参赛队伍内容确认\\
\bottomrule
\end{longtable}

\section{主要交互示例}

以下为归纳后的代表性提示，不是逐字完整聊天记录。

\begin{enumerate}
  \item “阅读2026年C题及全部附件，按四个子问题建立统一时间索引、能量平衡和储能约束；不得使用未来真实数据。”
  \item “对问题二设计只依赖过去日期的负荷/光伏预测和不确定性处理，说明冷启动、历史窗口、聚类退化和分位数定义。”
  \item “实现问题三0、6、12、18时更新，只调整尚未执行的合同；分别核算两种可能的合同结算解释。”
  \item “运行全年主策略和54组实验，输出题目指定Excel、代表日表、预测误差、物理残差、敏感性和更新成本阈值。”
  \item “生成覆盖四题的论文图，区分365日原始数据与334日策略统计，不把相关性写成提前知道价格。”
  \item “根据真实代码结果写LaTeX参考论文，摘要给出定量结果，明确2小时为合成更新、替代结算为同策略重计费，并检查PDF。”
\end{enumerate}

\section{AI建议的取舍与修改}

采纳的核心建议包括：用自然10分钟区间统一时标；储能SOC跨日连续；所有滚动预测只用已完成历史和已观测前缀；用联合残差刻画负荷与光伏误差，再以边际分位轨迹控制LP规模；对主结算与替代结算分别报告；将2小时结果表述为因果合成短临更新；在结论中加入每次额外更新的盈亏平衡成本。

未采纳或受限采纳的建议包括：不使用未来真实量进行完美信息优化；不把经验90\%区间称为概率保证；不把问题二到问题三的总差异全部归因于更新频率；不根据2025年同一批数据重新选择主参数并宣称全局最优；不把偏差修正描述为MAE改善；不声称四个代表日全部变好。

\section{待参赛队伍人工核验}

\begin{enumerate}
  \item 对照题目逐项确认时间标签、合同上调/下调费用和应急电价的解释。
  \item 阅读核心源代码，确认变量、约束、数据索引、历史窗口和输出表位置。
  \item 将论文中所有关键数值与五份结果工作簿、CSV汇总逐项核对。
  \item 核验参考文献与正文论述的匹配程度；如需更强理论论证，应阅读全文。
  \item 人工检查所有图表、公式、页码、摘要和附录，并按队伍真实思考改写文字。
  \item 根据竞赛平台最新要求完成文件命名、大小、上传、平台校验和最终提交，并留存成功凭证。
\end{enumerate}

\section{责任边界}

AI输出可能存在事实、计算、引用和表达错误。当前记录如实说明AI参与及尚未完成人工确认的状态。最终参赛作品的模型选择、结果真实性、规则合规性和提交责任均由参赛队伍承担。

\end{document}
